\subsection{Konfigurationsmodell}
\label{subsec:konfigurationsmodell}

Projektstruktur und Laufzeitkonfiguration sind getrennte Achsen:
\texttt{project-profile.toml} bestimmt die vorhandenen Features,
\path{config/environment.toml} den Variablenvertrag. Der zentrale Resolver
berücksichtigt nur Variablen der aktiven Features und wendet in absteigender
Priorität CLI-Override, Prozessumgebung, lokale \texttt{.env} und
Vertragsdefault an. Ableitungen wie die lokale API-URL bleiben als Quelle
nachvollziehbar \parencite{kleiveist2026configuration}.

Abbildung~\ref{fig:configuration-boundaries} zeigt die Adapter- und
Sicherheitsgrenzen. Vite bettet ausschließlich die explizit freigegebene
\texttt{VITE\_API\_BASE\_URL} ein. FastAPI verarbeitet serverseitige Werte
typisiert; \texttt{DATABASE\_URL} bleibt ein maskiertes Secret. Tooling gibt nur
relevante Werte an Kindprozesse weiter, und die Tauri-Umgebung entfernt bekannte
Secret-Namen. Ungültige Hosts, Ports, Origins oder erforderliche fehlende Werte
führen zu einem kontrollierten Abbruch, ohne den geheimen Eingabewert in der
Fehlermeldung auszugeben.

\begin{figure}[H]
  \centering
  \resizebox{0.95\textwidth}{!}{%
  \begin{tikzpicture}[node distance=6mm and 18mm]
    \node[casePrimary,minimum width=42mm] (cli) {CLI Override};
    \node[caseBox,minimum width=42mm,below=of cli] (env) {Process Environment};
    \node[caseBox,minimum width=42mm,below=of env] (dotenv) {.env};
    \node[caseBox,minimum width=42mm,below=of dotenv] (default) {Default};

    \draw[caseFlow] (cli) -- (env);
    \draw[caseFlow] (env) -- (dotenv);
    \draw[caseFlow] (dotenv) -- (default);

    \node[caseBox,minimum width=39mm,right=of cli] (frontend) {Frontend Values};
    \node[caseBox,minimum width=39mm,below=of frontend] (server) {Server-only Values};
    \node[caseOptional,minimum width=39mm,below=of server] (secrets) {Secrets};
    \begin{scope}[on background layer]
      \node[caseGroup,fit=(frontend)(server)(secrets),
        label={[font=\sffamily\bfseries\scriptsize,text=caseBlue]above:Sicherheitsgrenzen}] {};
    \end{scope}
  \end{tikzpicture}}
  \caption{Konfigurationspriorität und Sicherheitsgrenzen}
  \label{fig:configuration-boundaries}
\end{figure}

